\paragraph{What the study shows.} Under the inter-patient protocol, generative augmentation did not detectably improve the detection of rare classes over no augmentation in the pre-registered test (0 of 4 confirmatory comparisons rejected in its favour), and the same held with five seeds, with a patient-baseline RR feature, for a cnn made stronger on validation by patient-relative beats and with patient-transfer augmentation; on two external databases it was better than no augmentation in 1 and worse in 1 of 8 confirmatory comparisons. Whether a synthetic beat looks realistic and whether it helps a classifier on new patients are different questions, and this study answers the second.

\paragraph{Why.} The follow-up analyses P6 and P8 (Section~\ref{sec:posthoc}) point to the patients rather than the beats: the frozen cnn without augmentation reached S PR-AUC 0.942 on the training patients and 0.255 on the validation patients; generative augmentation narrowed the training-versus-validation gap in S PR-AUC in 4 of 4 cases (cnn cVAE: validation 0.255 to 0.375, training 0.942 to 0.949, cnn WGAN-GP: validation 0.255 to 0.313, training 0.942 to 0.938, rnn cVAE: validation 0.180 to 0.273, training 0.884 to 0.492, rnn WGAN-GP: validation 0.180 to 0.163, training 0.884 to 0.842), but none gave a detectably better test result; synthetic S beats were at least as concentrated on record 209 as the real ones (0.61 against 0.60). With so few training patients, a generator can add beats of those patients but not new ones; generators trained on more patients, or conditioned on the patient, are what this points to.

\paragraph{Absolute scores.} The test scores are low. Record 232 holds 0.75 of the test S beats and record 213 0.93 of the F beats, so the rare-class metrics largely measure two patients; record 208 holds 0.96 of the training F beats. A method that is good on average and poor on these patients scores low here, which is a property of the benchmark as much as of the method.

\paragraph{What the scores mean in the clinic.} A Holter reader needs two numbers for each class (Section~\ref{sec:medical}): how many of the patient's ectopic beats the software finds, and how many of its flags are right. On the test records the best pipeline (P18) found 0.91 of the S beats, and 0.50 of its S calls were S beats. On the 977 unused Icentia11k patients it found 0.59 of them, and 0.20 of its S calls were right, so about 8 of every 10 S flags would be false alarms for the reader to dismiss. It kept 0.96 of the normal beats as normal on the test records and 0.97 on the new patients: 27 to 38 of every 1,000 normal beats would be flagged as abnormal. Ventricular beats, whose wide QRS looks abnormal in almost every patient, carried over better than supraventricular ones: V F1 went from 0.936 on the test records to 0.648 on the new patients, S F1 from 0.647 to 0.294. Counting S beats, whose burden goes with later atrial fibrillation~\cite{binici2010}, would therefore still need a person to review every flag; at this scale the system is a screening aid to be checked, not a counter to be trusted.
